% Section 10 --- Related Work (~0.8 pages)
% Nine categories per v8.A-final outline.
%
% REVIEW STATUS: Draft. Categories 2 (arXiv 2410.11552), 3 (CCS 2025), 4
% (ethresear.ch) to be expanded with Task 1.5 annotated bibliography.
% All citation keys are placeholders pending references.bib population.

\section{Related Work}
\label{sec:related}

\paragraph{1. Classical Atomic Commitment.}
The two-phase commit protocol~\cite{gray1978notes} and its variants are the
canonical solution to distributed atomic commitment. Gray's original formulation
requires a coordinator with authority to abort partially-committed participants
--- precisely our capability (ii). Herlihy's universal construction~\cite{herlihy1991wait}
provides wait-free consensus but assumes a single shared memory object; the
multi-rollup setting has independent state machines with no shared write path.
The sagas pattern~\cite{garcia-molina1987sagas} allows long-running transactions
to commit partially, using compensating transactions for recovery. Sagas are
relevant to our analysis: a cross-rollup application using a saga pattern
implements capability (ii) for the non-final leg, and is therefore A2-compliant
within the saga's compensation scope. We note the structural similarity but
distinguish our contribution: we apply these classical notions to a setting
where the ``coordinator'' is a lazy sequencer with no execution access, which
classical atomic commitment does not address.

\paragraph{2. arXiv 2410.11552.}
% TODO (Task 1.5): fill with actual citation and summary.
% Placeholder: concurrent work on cross-rollup sequencing.
[CITATION] studies cross-rollup sequencing from a [FOCUS] perspective. Our
contribution is distinguished in [SPECIFIC WAY]: we provide the A2/A4 taxonomy
that separates execution outcome from economic externality, and demonstrate
the gap is non-vacuous via PoC attacks on deployed systems. [CITATION] does not
address [OUR SPECIFIC CONTRIBUTION].

\paragraph{3. CCS 2025 Denial of Sequencing.}
% TODO (Task 1.5): fill when paper is available.
[CITATION] identifies denial-of-sequencing as an attack vector on shared
sequencer architectures. Our contribution is orthogonal: we study the case
where the sequencer \emph{does} include the bundle (commits faithfully) but
the inclusion commitment still admits free-option attacks due to the execution
gap. Both attacks exploit properties of lazy sequencers; ours focuses on
selective exercise of commitments, not denial.

\paragraph{4. ethresear.ch Preconfirmation Taxonomies.}
Several ethresear.ch posts propose taxonomies and design criteria for
preconfirmations~\cite{Drake2023-BasedPreconf, Drake2023-BasedRollups,
Drake2023-ExecTickets}. These works focus primarily on L1 single-chain
preconfirmations (a validator committing to include a transaction in the
next block). Notably, the based-preconfirmations post~\cite{Drake2023-BasedPreconf}
defines three commitment levels (inclusion, execution state diffs, state root)
but makes no cross-rollup atomicity claims. Our contribution extends the
analysis to the cross-rollup setting, where the commitment spans two
independent STFs, and provides the first formal proof that A2 closure
requires capability (i) or (ii) regardless of bond size.

\paragraph{5. Unity is Strength.}
Obadia et al.~\cite{obadia2021unity} formalize cross-domain MEV and introduce
the notion of domain-crossing value extraction. Our A4 definition (economic
atomicity as a counterfactual payoff bound) extends this framework to the
sequencing commitment setting: the attacker's surplus is the option value
extracted from the commitment gap, and the domain boundary is the pair of
independent rollup STFs. The key distinction is that we study a structural
property of the commitment protocol, not just the value extraction outcome.

\paragraph{6. Flashbots Cross-Rollup Arbitrage Proposals.}
Flashbots has published proposals for cross-rollup MEV coordination and atomic
inclusion~\cite{flashbots2021mev, obadia2021unity}. These proposals identify
the coordination problem but propose economic solutions (bundler incentives,
relay-mediated ordering). Our Theorem~\ref{thm:a2-necessary} provides the
formal basis for why economic solutions (however well-designed) cannot close
the A2 gap: they are A4 mechanisms by construction.

\paragraph{7. EigenLayer Restaking.}
EigenLayer~\cite{eigenlayer-whitepaper} introduces restaking as a mechanism for
extending Ethereum's cryptoeconomic security to auxiliary services. Preconf
protocols using EigenLayer (e.g., Lagrange, Witness Chain) use restaked ETH as
collateral. As analyzed in Appendix~\ref{app:eigenlayer}, restaking increases
bond size ($B_{\max}$) without providing capability (i) or (ii). Our
Theorem~\ref{thm:bond-insufficient} shows that the resulting A4 guarantee
is structurally insufficient: the power-law option-value tail guarantees a
nonzero violation rate for any finite bond.

\paragraph{8. L1 MEV Literature.}
The MEV literature~\cite{daian2020flash, flashbots2021mev} establishes that
miners (and validators) extract value by reordering, censoring, or inserting
transactions. Our work is related but addresses a distinct mechanism: the
free option arises not from reordering within a single domain but from the
execution gap between a cross-domain commitment and execution. The sequencer
does not need to reorder transactions to exploit the option; they need only
observe partial execution outcomes and selectively honor or abandon the
commitment.

\paragraph{9. Smart Contract Vulnerability Analysis.}
The smart contract security literature~\cite{luu2016making} has
catalogued re-entrancy, oracle manipulation, and price impact attacks.
The PEFO attack we describe has structural similarity to price oracle
manipulation: the attacker creates conditions that cause one leg to succeed
and another to fail, extracting the price difference. The novel element is
that the attack surface is the sequencing commitment layer, not the individual
contract's logic. No single contract is buggy; the vulnerability is in the
interaction between the lazy sequencer and the independent-STF architecture.
